\begin{abstract}
We study cold detection from short speech recordings on the ComParE 2017 URTIC
corpus. The direction of the project was set by the properties of the dataset
rather than by a modeling idea. URTIC is cross-sectional: each of roughly 630
participants is recorded once, in a single health state; only 37 of the 210
participants in each official partition have a cold, and every participant
contributes many correlated chunks. Speaker identity is therefore perfectly
confounded with health state within the corpus, while the student release
provides no speaker identities with which to control for it. Prior work on the
challenge anticipated this failure mode, warning that such a classifier may
simply become a speaker recognizer, and reported severe development-to-Test
degradation~\cite{huckvale2017cold}. We consequently prioritize robustness and
proceed in three stages: first a trustworthy evaluation pipeline, then the
identification of feature signals that genuinely carry cold information, and
only then a comparison of model families.
We reconstruct pseudo-speaker groups from ECAPA embeddings and validate the
representation on two externally labeled corpora (ARI 0.996 and 0.963).
However, validating an embedding does not validate identities transferred with
it: a split reporting zero overlap under nearest-centroid assignments contains
201 of 210 crossing groups, and 98\% of its recordings belong to one, when
Development is clustered independently. Grouping construction alone moves a
fixed system by 0.076 UAR, twelve times its model-seed standard deviation. We
therefore confine model, fusion, and threshold selection to the fitting side,
and evaluate each frozen system once on the untouched official side.
Compact gain-invariant energy/pause descriptors, constant-Q summaries, and MFCC
path signatures provide complementary cold signal, whereas full
high-dimensional functional sets are weaker and more speaker-identifiable. Our
best hidden-Test system fuses seven gain-invariant energy/pause features with a
693-dimensional MFCC path log-signature and reaches 0.658 UAR, while two
substantially more complex HuBERT/eGeMAPS/CNN systems reach 0.656 and 0.657 at
markedly different operating points (accuracy 0.56 versus 0.77). Matched
negative controls are decisive for three de-confounding methods: a self-splice
control attributes the apparent augmentation effect to the splice operation
itself, a random-projection control reproduces nearly all of the speaker-probe
reduction attributed to contrastive learning, and per-epoch monitoring shows an
adversarial discriminator memorizing chunks rather than learning transferable
speaker structure. On a subject-limited corpus of this kind, evaluation design
and complementary signal determine what can be trusted far more than
architectural capacity does.
\end{abstract}
